% ---------------- SECCIÓN 3 ----------------
\section{Sistema de Evaluación y Rendimiento Académico}

\subsection{El Ciclo de 16 Semanas}
El ciclo universitario regular (que inicia el 23 de marzo) tiene una duración de 16 semanas. El sistema de evaluación clásico se divide de la siguiente manera:
\begin{itemize}
    \item \textbf{Semana 8 (Exámenes Parciales):} Evalúa todo lo visto desde la semana 1 hasta la 8.
    \item \textbf{Semana 16 (Exámenes Finales):} Evalúa todo lo visto desde la semana 8 hasta la 16.
    \item \textbf{Semana 17 (Sustitutorios):} El famoso "susti" reemplaza la nota más baja (ya sea del parcial o del final). Ojo, evalúa \textbf{todo el curso} (semana 1 a la 16) y es de doble filo si sacas menos nota de la que ya tenías.
\end{itemize}

\subsection{Consecuencias de desaprobar (Bikas y Trikas)}
Eviten jalar cursos a toda costa. El reglamento de la UNI es estricto:
\begin{itemize}
    \item \textbf{Bika (Jalar 2 veces un curso):} Pasas a ser "alumno en riesgo académico". Entras a tutoría obligatoria, recibes atención psicológica y, lo peor, se limita tu matrícula a menos de 4 cursos, lo que te atrasa al menos un año entero.
    \item \textbf{Trika (Jalar 3 veces un curso):} "Alumno en alto riesgo". Implica separación automática de la universidad por un año. Al volver, \textbf{solo} podrás matricularte en ese curso.
    \item \textbf{Cuarta vez:} Si al volver del baneo vuelves a desaprobar ese curso, la expulsión de la UNI es permanente.
\end{itemize}
